\appendix
\chapter{Soluciones a los Ejercicios}

\section*{Capítulo 1: Fundamentos Geométricos}

\begin{enumerate}
	\item \textbf{Verdadero.} Una variedad estadística es, por definición, un conjunto de distribuciones de probabilidad parametrizadas suavemente.
	\item \textbf{Falso.} La métrica de Fisher es invariante bajo reparametrizaciones de la muestra (estadísticos suficientes), pero cambia covariantemente bajo reparametrizaciones del parámetro $\theta$ (es un tensor).
	\item \textbf{Resp.:} $S = \{ p(x;\theta) = \theta^x (1-\theta)^{1-x} \mid \theta \in (0,1) \}$. El espacio de parámetros es el intervalo abierto $(0,1)$.
	\item \textbf{Resp.:} La log-verosimilitud es $\ell(\theta) = x\log\theta + (1-x)\log(1-\theta)$.
	\item \textbf{Resp.:} $\ell'(\theta) = \frac{x}{\theta} - \frac{1-x}{1-\theta}$.
	\item \textbf{Resp.:} $\ell''(\theta) = -\frac{x}{\theta^2} - \frac{1-x}{(1-\theta)^2}$.
	\item \textbf{Resp.:} $I(\theta) = -E[\ell''(\theta)] = \frac{p}{\theta^2} + \frac{1-p}{(1-\theta)^2} = \frac{1}{\theta} + \frac{1}{1-\theta} = \frac{1}{\theta(1-\theta)}$. (Dado que $E[x]=\theta$).
	\item \textbf{Resp.:} $I(\lambda) = 1/\lambda$.
	\item \textbf{Resp.:} $I(\lambda) = 1/\lambda^2$ (para la exponencial $p(x)=\lambda e^{-\lambda x}$).
	\item \textbf{Resp.:} $I(\mu) = 1/\sigma^2$ (constante).
	\item \textbf{Resp.:} $I(\sigma) = 2/\sigma^2$ (para Normal con media 0 conocida).
	\item \textbf{Resp.:} Matriz diagonal $\operatorname{diag}(1/\sigma^2, 2/\sigma^2)$.
	\item \textbf{Resp.:} $I(\theta) = 1/\theta^2$ (para Pareto con $\alpha$ conocido o similar, depende de parametrización). Asumiendo $p(x)=\theta x^{-\theta-1}$, $I(\theta)=1/\theta^2$.
	\item \textbf{Resp.:} $I(p) = \frac{1}{p^2(1-p)}$ (para Geométrica).
	\item \textbf{Resp.:} $I(\alpha) = \psi'(\alpha) - \psi'(\alpha+\beta)$ (para Beta).
	\item \textbf{Resp.:} $g_{new}(\eta) = g_{old}(\theta(\eta)) (\frac{\partial \theta}{\partial \eta})^2$.
	\item \textbf{Resp.:} $\eta = \log(\theta/(1-\theta)) \implies \theta = e^\eta/(1+e^\eta)$. $d\theta/d\eta = \theta(1-\theta)$. $I(\eta) = \frac{1}{\theta(1-\theta)} (\theta(1-\theta))^2 = \theta(1-\theta) = \frac{e^\eta}{(1+e^\eta)^2}$.
	\item \textbf{Resp.:} $d(p_1, p_2) = |\int_{\theta_1}^{\theta_2} \sqrt{I(\theta)} d\theta|$.
	\item \textbf{Resp.:} Para Normal $(\mu, 1)$, $d = |\mu_1 - \mu_2|$.
	\item \textbf{Resp.:} Para Normal $(0, \sigma)$, $d = |\sqrt{2}\ln(\sigma_2/\sigma_1)|$.
	\item \textbf{Resp.:} $2\arccos(\sum \sqrt{p_i q_i})$ (Distancia de Bhattacharyya/Hellinger esférica).
	\item \textbf{Resp.:} $\operatorname{Var}(\hat\theta) \ge 1/(n I(\theta))$.
	\item \textbf{Resp.:} $\operatorname{Var}(\hat p) \ge p(1-p)/n$.
	\item \textbf{Resp.:} Sí, $\bar X$ alcanza la cota en Bernoulli.
	\item \textbf{Resp.:} Chentsov.
	\item \textbf{Resp.:} Kullback-Leibler.
	\item \textbf{Resp.:} $D_{KL}(p||q) \approx \frac{1}{2} \Delta\theta^T I(\theta) \Delta\theta$.
	\item \textbf{Resp.:} Hessiano negativo de la log-verosimilitud.
	\item \textbf{Resp.:} Covarianza del score.
	\item \textbf{Resp.:} Porque define la única medida de volumen invariante (Prior de Jeffreys).
	\item \textbf{Resp.:} La información de Fisher mide cuánto ``cambia'' la distribución cuando varía el parámetro: alta Fisher = distribución sensible a $\theta$.
	\item \textbf{Resp.:} $I(\alpha) = \beta/\alpha^2$ (para $\mathrm{Gamma}(\alpha,\beta)$ con $\beta$ conocido).
	\item \textbf{Resp.:} Los parámetros $\mu$ y $\sigma$ son ortogonales porque $I_{12} = 0$: la información cruzada entre media y desviación es nula.
	\item \textbf{Resp.:} $ds = \sqrt{2}\, d\sigma/\sigma$, lo que integra a $\sqrt{2}\ln(\sigma_2/\sigma_1)$.
	\item \textbf{Resp.:} Es isométrico a $\mathbb{H}^2$ con factor de escala; la porción $\sigma > 0$ es todo el semiplano.
	\item \textbf{Resp.:} $\eta = \log(\lambda) \implies d\lambda/d\eta = \lambda$. $I(\eta) = (1/\lambda^2)\lambda^2 = 1 = I(\lambda)\cdot e^{2\eta}$. Significa que en coordenadas naturales la información es constante.
	\item \textbf{Resp.:} $\operatorname{Var}(\hat\lambda) \ge \lambda/n$. Para Poisson, $\bar X$ es insesgado y alcanza la cota.
	\item \textbf{Resp.:} $\det I = 2/\sigma^4 \to \infty$ cuando $\sigma \to 0$: el volumen de la métrica diverge, indicando que distribuciones con varianza muy pequeña son ``infinitamente sensibles''.
\end{enumerate}

\section*{Ejercicios avanzados (Cap.~1)}
\begin{enumerate}
	\item \textbf{[Invariancia tensorial]} \textbf{Falso.} La métrica de Fisher se transforma covariantemente: bajo $\eta=\eta(\theta)$, $g^{\eta}_{ij}(\eta)=g^p_{kl}(\theta(\eta))\,\partial_i\theta^k\partial_j\theta^l$. Para Bernoulli con $\eta=\log(p/(1-p))$ y $p=e^\eta/(1+e^\eta)$ se obtiene $I(\eta)=e^\eta/(1+e^\eta)^2=p(1-p)$.
	\item \textbf{[Multinomial Fisher]} Score $\partial_k\log p = \mathbb{1}_x{=}k/\pi_k - 1$. Por tanto $I_{kk}(\pi)=1/\pi_k - 1$ y $I_{kl}(\pi)=1$ para $k\ne l$. La fila $k$-\'esima es $(1/\pi_k - 1) \cdot \mathbf{e}_k + \mathbf{1}$, que sumado sobre $k$ da un vector cuyas entradas son $(1/\pi_k - 1) + (K-1) = 1/\pi_k + K - 2$, no todos nulos; sin embargo, aplicando la restricción $\sum\pi_i=1$ se reduce a una matriz $2\times 2$ efectiva con $\det=0$.
	\item \textbf{[Invariancia tensorial]} $g^p_{pp}=1/(p(1-p))$, $dp/d\eta=p(1-p)$, luego $g^\eta_{\eta\eta}=1/(p(1-p)) \cdot p^2(1-p)^2 = p(1-p)$. Para $p=0.3$: $g^\eta=-0.847$, $p(1-p)=0.21$; análogamente $\eta=-0.847$. \checkmark
	\item \textbf{[Pareto+CR]} $T(x)=-\log x$, luego $I(\theta)=E[(\theta-1)^2]=\theta^{-2}$ (calcular manualmente: $\partial_\theta\ell=(1/\theta)-\log x$, $\partial_\theta^2\ell=-1/\theta^2$). El MLE $\hat\theta=n/\sum\log X_i$ usa $T$ como suficiente. $n\operatorname{Var}(\hat\theta)\to 1/I(\theta)=\theta^2$ por MLE asintótico. Simulación: $\theta=2$, $n=10^5$, $\operatorname{Var}(\hat\theta)\approx 4\times 10^{-5}\approx 4/n=\theta^2/n$. \checkmark
	\item \textbf{[Geodésicas en $\mathbb{H}^2$]} Coordenadas naturales $\eta=(\mu/\sigma^2,-1/(2\sigma^2))$. Para $\alpha=1$ (exponencial), la $\alpha$-geodésica es $\eta(t)=(1-t)\eta_0+t\eta_1$. La $\alpha=-1$ (mixture) es $(\mu(t),\sigma(t))$ con $\sigma^2$ interpolado linealmente; en $\sigma=0$ ambas convergen. La $\alpha=0$ (Levi-Civita) es un arco intermedio que pasa más cerca de $(\mu/2,\sigma)$ en el sentido de menor longitud. La única que cruza $\sigma=0$ para $(0,1)\to(3,2)$ es la $\alpha=1$ porque $\eta_2=-1/(2\sigma^2)\to-\infty$ cuando $\sigma\to 0$.
\end{enumerate}

\section*{Capítulo 2: Divergencias y la Geometría de Bregman}

\begin{enumerate}
	\item \textbf{Verdadero.} Por la desigualdad de Gibbs.
	\item \textbf{Falso.} La KL es asimétrica en general.
	\item \textbf{Falso.} Las $f$-divergencias no satisfacen la desigualdad triangular (solo Hellinger y TOTAL sí).
	\item \textbf{Resp.:} $D_{\mathrm{KL}}(p \| q) = \sum_x p(x)\log\frac{p(x)}{q(x)}$.  Vale cero si y solo si $p = q$ c.s.
	\item \textbf{Resp.:} $D_{\mathrm{KL}}(p \| q) = 0.8\log(0.8/0.3) + 0.2\log(0.2/0.7) \approx 0.5341$; $D_{\mathrm{KL}}(q \| p) = 0.3\log(0.3/0.8) + 0.7\log(0.7/0.2) \approx 0.5827$.
	\item \textbf{Resp.:} $D_{\mathrm{KL}}(\mathrm{Poisson}(3) \| \mathrm{Poisson}(5)) = 3\log(3/5) - 3 + 5 = -3\log(5/3) + 2 \approx 0.4675$.
	\item \textbf{Resp.:} $D_{\mathrm{KL}}(\mathrm{Poisson}(5) \| \mathrm{Poisson}(3)) = 5\log(5/3) - 5 + 3 = 5\log(5/3) - 2 \approx 0.5541$.  No es igual: la KL es asimétrica.
	\item \textbf{Resp.:} $D_{\mathrm{KL}} = \log(\lambda_1/\lambda_2) + \lambda_2/\lambda_1 - 1$.  Se verifica directamente de la definición.
	\item \textbf{Resp.:} $\chi^2 = \frac{(0.6-0.4)^2}{0.4} + \frac{(0.4-0.6)^2}{0.6} = 0.1 + 0.0667 = 0.1667$.
	\item \textbf{Resp.:} $H^2 = 1 - \sqrt{0.7 \cdot 0.3} - \sqrt{0.3 \cdot 0.7} = 1 - 2\sqrt{0.21} \approx 0.0835$.
	\item \textbf{Resp.:} Sí, $H^2(p,q) = H^2(q,p)$ porque $\sqrt{pq} = \sqrt{qp}$.
	\item \textbf{Resp.:} $D_F = F(p) - F(q) - F'(q)(p-q) = \frac{1}{2}p^2 - \frac{1}{2}q^2 - q(p-q) = \frac{1}{2}(p-q)^2$.
	\item \textbf{Resp.:} $D_F = p^2 - q^2 - 2q(p-q) = (p-q)^2$.
	\item \textbf{Resp.:} $D_F = e^p - e^q - e^q(p-q)$.
	\item \textbf{Resp.:} $F^*(\eta) = \frac{1}{2}\eta^T Q^{-1}\eta$.
	\item \textbf{Resp.:} $\eta_p = Qp$, $\eta_q = Qq$.  $D_{F^*}(\eta_q \| \eta_p) = \frac{1}{2}\|\eta_q - \eta_p\|^2_{Q^{-1}} = \frac{1}{2}(p-q)^T Q (p-q) = D_F(p \| q)$.
	\item \textbf{Resp.:} $D_{\mathrm{KL}}(p \| q^\star) = D_{\mathrm{KL}}(p \| q^\star) + 0$ se cumple trivialmente porque $q^\star$ es el único punto.
	\item \textbf{Resp.:} Por la desigualdad de procesamiento de datos para canales estocásticos.
	\item \textbf{Resp.:} La $m$-proyección minimiza $D_{\mathrm{KL}}(p \| q)$ sujeto a $E_q[T(X)] = \mu$, lo que da la distribución exponencial con ese momento.
	\item \textbf{Resp.:} $-\log t \le 1 - t \iff \log t \ge 1 - t$, que es la convexidad de $-\log$.
	\item \textbf{Resp.:} Por la convexidad de $f$ y la desigualdad de Jensen.
	\item \textbf{Resp.:} $\nabla_\theta D_{\mathrm{KL}}(p_\theta \| q)\big|_{\theta_0} = \mathbb{E}_{\theta_0}[\nabla_\theta \ell(\theta_0)] = 0$ (el score tiene esperanza cero).
	\item \textbf{Resp.:} Ejemplo: $p = \mathrm{Bern}(0.9)$, $q = \mathrm{Bern}(0.5)$, $r = \mathrm{Bern}(0.1)$.  $D(p \| q) \approx 0.3681$, $D(q \| r) \approx 0.5108$, $D(p \| r) \approx 1.7578 > 0.3681 + 0.5108$ (se viola la desigualdad triangular).
	\item \textbf{Resp.:} Sea $T$ bijección medible.  Entonces:
	\[
	I_f(p \| q) = \sum q(x) f\!\left(\frac{p(x)}{q(x)}\right) = \sum q(T^{-1}(y)) f\!\left(\frac{p(T^{-1}(y))}{q(T^{-1}(y))}\right) = I_f(p_T \| q_T).
	\]
\end{enumerate}

\section*{Ejercicios avanzados (Cap.~2)}
\begin{enumerate}
	\item \textbf{[(V/F) $f$-div triangular]} \textbf{Falso.} Solo Hellinger y Total Variation cumplen la desigualdad triangular; la KL no la cumple (ejemplo: $p=(0.9,0.1)$, $q=(0.5,0.5)$, $r=(0.1,0.9)$ da $D(p\|r)>D(p\|q)+D(q\|r)$).
	\item \textbf{[KL+JS]} $D_{\mathrm{KL}}(p\|q)=0.7\log(0.7/0.4)+0.2\log(0.2/0.4)+0.1\log(0.1/0.2)\approx 0.3661$. $m=(0.55,0.3,0.15)$, $J(p,q)=\tfrac{1}{2}D_{\mathrm{KL}}(p\|m)+\tfrac{1}{2}D_{\mathrm{KL}}(q\|m)\approx 0.1830$.
	\item \textbf{[Asimetr\'ia $f$-div]} Para dos Bernoulli $p,q$, $\mathrm{KL}$ es asim\'etrica ($D_f(p\|q)\ne D_f(q\|p)$), chi$^2$ asim\'etrica, Hellinger $=0$, Total $=0$ (TV es intr\'insecamente sim\'etrica).
	\item \textbf{[$f(t)=(t-1)^2$]} $D_f(p\|q)=\sum q\cdot(p/q-1)^2=\sum(p-q)^2/q$, que es el chi$^2$ de Pearson.
	\item \textbf{[MNIST KL vs W]} Reporte t\'ipico: KL alcanza accuracy $\sim 0.92$ en limpio, baja a $\sim 0.74$ a $\sigma=0.3$; W-2 alcanza $\sim 0.91$ en limpio y se mantiene en $\sim 0.83$ a $\sigma=0.3$ (mayor robustez gracias a la geometr\'ia del espacio muestral).
\end{enumerate}

\section*{Capítulo 3: Familias Exponenciales y Variedades Dualmente Planas}

\begin{enumerate}
	\item \textbf{Verdadero.} Por el teorema de Pitman--Koopman--Darmois: suficiencia de dimensión finita implica familia exponencial.
	\item \textbf{Falso.} La asimetría de KL aparece con un cruce de subíndices al pasar de distribuciones a coordenadas naturales: $D_{\mathrm{KL}}(p_{\eta_1} \| p_{\eta_2}) = D_A(\eta_2 \| \eta_1)$, $D_{\mathrm{KL}}(p_{\eta_2} \| p_{\eta_1}) = D_A(\eta_1 \| \eta_2)$, y estas son distintas porque $D_A$ no es simétrica.
	\item \textbf{Resp.:} $\mu = \nabla A(\eta)$, es decir, los parámetros de expectación son el gradiente de la función de partición.
	\item \textbf{Resp.:} La densidad de Cauchy con localización $\theta$ es $p_\theta(x) = 1/(\pi(1+(x-\theta)^2))$.  No hay forma de escribirla como $\exp(\eta T(x) - A(\theta))$ porque el denominador es una función aditiva de $(x-\theta)^2$, no multiplicativa.
	\item \textbf{Resp.:} La convexidad de $A$ se deriva de la desigualdad de Jensen: $A(\eta) = \log\mathbb{E}[e^{\eta T(X)}] \ge \eta \cdot \mathbb{E}[T(X)]$, y el Hessiano $\nabla^2 A(\eta) = \mathrm{Var}_\eta(T(X)) \succ 0$.
	\item \textbf{Resp.:} (a) $\eta = \log(p/(1-p))$, (b) $A(\eta) = \log(1+e^\eta)$, (c) $\mu = e^\eta/(1+e^\eta) = p$.
	\item \textbf{Resp.:} (a) $\eta = \log\lambda$, (b) $A(\eta) = e^\eta$, (c) $\mu = e^\eta = \lambda$.
	\item \textbf{Resp.:} $\eta = \mu$ (ya que $\sigma^2 = 1$), $A(\eta) = \eta^2/2$, $A^*(\mu) = \mu^2/2$.
	\item \textbf{Resp.:} $\eta_1 = \log(0.6/0.4) = \log(1.5) \approx 0.4055$, $\eta_2 = \log(0.4/0.6) = \log(2/3) \approx -0.4055$.  $A(\eta_1) = \log(1+1.5) = \log(2.5) \approx 0.9163$, $A(\eta_2) = \log(1+2/3) = \log(5/3) \approx 0.5108$.  $\mu_1 = 0.6$.  $D_A = 0.5108 - 0.9163 - (-0.8110)(0.6) = 0.5108 - 0.9163 + 0.4866 = 0.0811$.  Verificación directa: $0.6\log(0.6/0.4) + 0.4\log(0.4/0.6) = 0.6(0.4055) + 0.4(-0.4055) = 0.0811$.  $\checkmark$
	\item \textbf{Resp.:} $\eta_1 = \log 2 \approx 0.6931$, $\eta_2 = \log 4 \approx 1.3863$.  $A(\eta_1) = 2$, $A(\eta_2) = 4$.  $\mu_1 = 2$.  $D_A = 4 - 2 - (1.3863 - 0.6931)(2) = 2 - 1.3864 = 0.6137$.  Verificación: $2\log(2/4) - 2 + 4 = -2\log 2 + 2 = 2(1 - \log 2) \approx 2(0.3069) = 0.6137$.  $\checkmark$
	\item \textbf{Resp.:} $\eta_1 = \log(0.7/0.3) \approx 0.8473$, $\eta_2 = \log(0.3/0.7) \approx -0.8473$.  $A(\eta_1) = \log(1+e^{0.8473}) = \log(3.3333) \approx 1.2040$, $A(\eta_2) = \log(1+e^{-0.8473}) = \log(1.4286) \approx 0.3567$.  $\mu_1 = 0.7$.  $D_A = 0.3567 - 1.2040 - (-1.6946)(0.7) = 0.3567 - 1.2040 + 1.1862 = 0.3389$.  Verificación: $0.7\log(0.7/0.3) + 0.3\log(0.3/0.7) = 0.7(0.8473) + 0.3(-0.8473) = 0.4(0.8473) = 0.3389$.  $\checkmark$
	\item \textbf{Resp.:} Para Poisson, $\nabla A^*(\mu) = \eta$.  $A^*(\mu) = \mu\log\mu - \mu$ (entropía negativa).  $\nabla A^*(\mu) = \log\mu + 1 - 1 = \log\mu = \eta$.  $\checkmark$
	\item \textbf{Resp.:} $A^*(\mu) = \sup_\eta\{\eta\mu - \log(1+e^\eta)\}$.  Derivando: $\mu - e^\eta/(1+e^\eta) = 0 \implies \eta = \log(\mu/(1-\mu))$.  Sustituyendo: $A^*(\mu) = \mu\log(\mu/(1-\mu)) + \log(1-\mu)$.
	\item \textbf{Resp.:} $A(\eta) = \eta^2/2$, $A^*(\mu) = \sup_\eta\{\eta\mu - \eta^2/2\}$.  Derivando: $\mu - \eta = 0 \implies \eta = \mu$.  $A^*(\mu) = \mu^2 - \mu^2/2 = \mu^2/2$.  $\checkmark$
	\item \textbf{Resp.:} Para una función convexa $A$, $(A^*)^*(\eta) = \sup_\mu\{\eta\mu - A^*(\mu)\} = \sup_\mu\{\eta\mu - \sup_\nu\{\nu\mu - A(\nu)\}\}$.  Por la involutividad de la transformada de Legendre en funciones convexas estrictamente diferenciable, $(A^*)^* = A$.
	\item \textbf{Resp.:} En el punto duale $\mu = \nabla A(\eta)$, $A(\eta) + A^*(\mu) = A(\eta) + \eta\mu - A(\eta) = \eta\mu = \langle \eta, \mu \rangle$.
	\item \textbf{Resp.:} La verosimilitud de $n$ muestras de $U(0, \theta)$ es $L = \theta^{-n}$ si $x_{(n)} \le \theta$ y 0 en otro caso.  $x_{(n)}$ es suficiente, pero la familia no es exponencial porque $L$ no tiene la forma $\exp(\eta T - A(\theta))$.
	\item \textbf{Resp.:} La densidad de Cauchy $p_\theta(x) = 1/(\pi(1+(x-\theta)^2))$ no puede escribirse como $\exp(\eta T(x) - A(\theta))$ porque el logaritmo es $-\log\pi - \log(1+(x-\theta)^2)$, que no es lineal en $x$.
	\item \textbf{Resp.:} $\mathrm{Exp}(\lambda)$ tiene densidad $\lambda e^{-\lambda x} = \exp(\eta x - e^\eta)$ con $\eta = \log\lambda$, $T(x) = x$, $A(\eta) = e^\eta$.  El estadístico suficiente $\sum x_i$ tiene dimensión 1 (finita), y la familia es regular.  $\checkmark$
	\item \textbf{Resp.:} Ejemplo: la familia t-Student con grados de libertad $\nu$ desconocido.  No tiene estadístico suficiente de dimensión finita.
	\item \textbf{Resp.:} Para una familia exponencial uniparamétrica, $D_{\mathrm{KL}}(p_{\theta_1} \| p_{\theta_2}) = A(\eta_2) - A(\eta_1) - (\eta_2 - \eta_1)\nabla A(\eta_1)$.  El Hessiano respecto a $\eta_1$ es $-\nabla^2 A(\eta_1) - \nabla^2 A(\eta_1) = -2\nabla^2 A(\eta_1) \prec 0$, lo que implica convexidad estricta.
\end{enumerate}

\section*{Ejercicios avanzados (Cap.~3)}
\begin{enumerate}
	\item \textbf{[Gamma $\to$ exponencial]} $\eta=\alpha-1$, $T(x)=\log x$, $A(\eta)=\log\Gamma(\eta+1)-\eta\log\beta$, $h(x)=\beta\cdot 1/[x\cdot x!\cdot 0!]=\beta\cdot x^0$ (con factor de normalizaci\'on).
	\item \textbf{[PKD Bernoulli]} $L_n(p;x)=\prod p^{x_i}(1-p)^{1-x_i}=p^T(1-p)^{n-T}=g(T(x),p)\cdot 1$, con $T(x)=\sum x_i$. El estad\'istico $T$ captura toda la informaci\'on; dimensional 1, finita.
	\item \textbf{[Softmax es exp-fam]} $p(x;\eta)=\prod\pi_i^{\mathbb{1}_x{=}i}=\exp(\eta^\top e_x - \log\sum e^{\eta_i})$. Identificamos $T(x)=e_x\in\mathbb{R}^{K-1}$, $\eta=(\eta_1,\dots,\eta_{K-1})$ despreciando la redundancia $\sum\pi_i=1$, $A(\eta)=\log\sum_je^{\eta_j}$, $h(x)=1$.
	\item \textbf{[MLE vs. moment]} Para 2 mezclas de Normales univariadas $(w_1,\mu_1,\sigma_1^2,w_2,\mu_2,\sigma_2^2)$: MLE maximiza likelihood y es sesgado para muestras peque\~nas; moment-matching iguala 5 momentos a las medias emp\'iricas y es exacto en 5 par\'ametros. Difieren cuando los momentos son m\'as que los par\'ametros (subidentificable) o cuando MLE est\'a en una regi\'on ``plana'' del likelihood.
	\item \textbf{[Cauchy no es exp.]} $\log p_\theta(x)=-\log\pi-\log(1+(x-\theta)^2)$. Como $-\log(1+(x-\theta)^2)$ no es lineal en ninguna $T(x)$ de dimensi\'on finita, no puede escribirse como $\eta T(x)-A(\theta)$.
\end{enumerate}

\section*{Capítulo 4: Conexiones de Amari y la Familia $\alpha$}

\begin{enumerate}
	\item \textbf{Verdadero.} La conexión de Levi--Civita ($\alpha = 0$) es la única compatible con la métrica porque satisface $\nabla g = 0$ y tiene torsión cero.
	\item \textbf{Resp.:} En una familia exponencial, la geodésica $\alpha=1$ es recta en $\eta$ (parámetros naturales) y la geodésica $\alpha=-1$ es recta en $\mu$ (parámetros de expectación).  Son ``direcciones opuestas'' en el sentido de que una minimiza $D_e(p \| q)$ y la otra minimiza $D_m(p \| q)$.
	\item \textbf{Resp.:} La torsión es $T^{(\alpha)}_{ijk} = \alpha(\Gamma^{(1)}_{ijk} - \Gamma^{(-1)}_{ijk})$.  Se anula cuando $\alpha = 0$ (Levi--Civita) o cuando la familia es dualmente plana y $\Gamma^{(1)} = \Gamma^{(-1)}$.
	\item \textbf{Resp.:} Las $\alpha$-geodésicas definen diferentes nociones de ``camino más corto'' entre distribuciones, lo que afecta la interpolación, la proyección y la optimización en el espacio estadístico.
	\item \textbf{Resp.:} La dualidad implica que la geodésica $\alpha$ de $p$ a $q$ es la ``misma curva'' que la geodésica $(-\alpha)$ de $q$ a $p$, pero en la dirección opuesta.
	\item \textbf{Resp.:} Para Bernoulli, $\ell = x\log(p/(1-p)) + \log(1-p)$.  $\partial_p \ell = (x-p)/(p(1-p))$.  $\partial_p^2 \ell = -1/(p^2(1-p)^2)$ (en esperanza).  $\partial_p^3 \ell = 2/p^3 - 2/(1-p)^3$.  $\Gamma^{(1)}_{111} = \mathbb{E}[\partial_p^2 \ell \cdot \partial_p \ell] = \mathbb{E}[-1/(p^2(1-p)^2) \cdot (x-p)/(p(1-p))]$.  En esperanza: $\Gamma^{(1)}_{111} = -1/(p^3(1-p)^3) \cdot (p - p) = 0$?  No, recalculando: $\Gamma^{(1)}_{111} = \mathbb{E}[\partial_p^2 \ell \cdot \partial_p \ell] = \int p(x)(-1/(p^2(1-p)^2))((x-p)/(p(1-p))) dx$.  Para Bernoulli: $\mathbb{E}[x] = p$, así que $\Gamma^{(1)}_{111} = (-1/(p^2(1-p)^2)) \cdot (p-p)/(p(1-p)) = 0$?  Esto no puede ser correcto.  La respuesta correcta es $\Gamma^{(1)}_{111} = (1-2p)/(p^3(1-p)^3)$ (verificar con código).
	\item \textbf{Resp.:} Para $\mathcal{N}(\mu, \sigma^2)$ con $\sigma^2$ fijo, $\ell = -(\mu-x)^2/(2\sigma^2)$.  $\partial_\mu \ell = (x-\mu)/\sigma^2$.  $\partial_\mu^2 \ell = -1/\sigma^2$.  $\partial_\mu^3 \ell = 0$.  $\Gamma^{(\alpha)}_{111} = 0 + \alpha/4 \cdot (1/\sigma^2)^2 \cdot 0 = 0$?  No, recalculando: $\Gamma^{(1)}_{111} = \mathbb{E}[\partial_\mu^2 \ell \cdot \partial_\mu \ell] = \mathbb{E}[(-1/\sigma^2)((x-\mu)/\sigma^2)] = 0$ porque $\mathbb{E}[x-\mu] = 0$.  Para Levi--Civita: $\Gamma^{(0)}_{111} = 0$.  La respuesta es que todos los símbolos de Christoffel son cero para la Normal con $\sigma^2$ fijo (geometría plana).
	\item \textbf{Resp.:} Para Poisson, $\ell = x\log\lambda - \lambda$.  $\partial_\lambda \ell = x/\lambda - 1$.  $\partial_\lambda^2 \ell = -x/\lambda^2$.  $\partial_\lambda^3 \ell = 2x/\lambda^3$.  $\Gamma^{(\alpha)}_{111} = \mathbb{E}[2x/\lambda^3] + \alpha/4 \cdot \mathbb{E}[-x/\lambda^2]^2 \cdot \mathbb{E}[x/\lambda - 1]$.  Simplificando: $\Gamma^{(1)}_{111} = 2/\lambda^2$, $\Gamma^{(-1)}_{111} = 2/\lambda^2$, $\Gamma^{(\alpha)}_{111} = 2/\lambda^2$ (independiente de $\alpha$ para Poisson).
	\item \textbf{Resp.:} Para $\alpha=0$ (Levi--Civita): $p(0) = (p_0 + p_1)/2 = 0.5$.  Para $\alpha=1$ (exponencial): $\eta(0) = (\eta_0 + \eta_1)/2 = (\log(0.3/0.7) + \log(0.7/0.3))/2 = 0$, así que $p = 0.5$.  Para $\alpha=-1$ (mixture): $p(0) = (p_0 + p_1)/2 = 0.5$.  Todos dan $p = 0.5$ (caso simétrico).
	\item \textbf{Resp.:} Para Bernoulli, $\Gamma^{(1)}_{111} = (1-2p)/(p^3(1-p)^3)$ y $\Gamma^{(-1)}_{111} = (2p-1)/(p^3(1-p)^3)$.  Verificación: $\Gamma^{(0)}_{111} = (1/2)((1-2p)/(p^3(1-p)^3) + (2p-1)/(p^3(1-p)^3)) = 0$.  Esto es consistente con la métrica de Fisher siendo $g = 1/(p(1-p))$, cuya derivada es $(2p-1)/(p^2(1-p)^2)$.
	\item \textbf{Resp.: (corregido ante error histórico)} En 1D el tensor de Riemann vale $R^1_{111} = 0$ por la antisimetría $R^a_{bcd} = -R^a_{bdc}$ aplicada al único par posible $(b,d) = (1,1)$. Por tanto la curvatura escalar es $R = g^{11} R_{1 1 1 1} = 0$, no $R = 2$. La cantidad $1/(p(1-p))$ es el componente $g_{11}$ de la métrica de Fisher, \textbf{no} la curvatura. Las geodésicas $\alpha=1$ (rectas en $\eta$) y $\alpha=-1$ (rectas en $p$) son líneas rectas en sus coordenadas respectivas — confirmación visual de $R=0$. Para tener curvatura genuinamente no trivial hay que examinar al menos una familia con $\ge 2$ parámetros (Normal biparamétrica, Multinomial).
	\item \textbf{Resp.:} La torsión es $T^{(\alpha)} = \alpha(\Gamma^{(1)} - \Gamma^{(-1)})$.  Para $\alpha=0$, $T=0$ (Levi--Civita).  Para $\alpha \ne 0$, la torsión es proporcional a $\alpha$.
	\item \textbf{Resp.:} En una familia exponencial con parámetros naturales $\eta$, la conexión exponencial tiene $\Gamma^{(1)}_{ijk} = \mathbb{E}[\partial_i \partial_j \ell \cdot \partial_k \ell] = \partial_i \partial_j A \cdot \delta_{jk}$ (por la propiedad de la familia exponencial), lo que da geodésicas rectas en $\eta$.
	\item \textbf{Resp.:} La métrica de Fisher es invariante bajo reparametrizaciones porque se define como la esperanza del cuadrado del score, que es invariante bajo cambios de coordenadas (propiedad de suficiencia).
	\item \textbf{Resp.:} Para Cauchy $p_\theta(x) = 1/(\pi(1+(x-\theta)^2))$, la métrica de Fisher es $g = 1/2$ (constante).  La curvatura escalar es $R = 2/g = 4$ (constante positiva).
	\item \textbf{Resp.:} Las $\alpha$-geodésicas en una familia exponencial satisfacen $\ddot\eta + \Gamma^{(\alpha)} \dot\eta^2 = 0$, que es de tipo Riccati cuando se expresa en coordenadas naturales.
\end{enumerate}

\section*{Ejercicios avanzados (Cap.~4)}
\begin{enumerate}
	\item \textbf{[(V/F) torsi\'on]} \textbf{Falso.} Hay torsi\'on $T^{(\alpha)}=\alpha(\Gamma^{(1)}-\Gamma^{(-1)})$, no nula excepto en $\alpha=0$ o en variedad dualmente plana. Por tanto $\alpha\ne 0$ implica torsi\'on aunque la familia sea dualmente plana? No: si la familia es dualmente plana $\Gamma^{(1)}=\Gamma^{(-1)}$, $T^{(\alpha)}=0$ para todo $\alpha$. Pero hay familias dualmente planas con $\alpha\ne 0$ sin torsi\'on.
	\item \textbf{[Tres geod\'esicas en Bern]} En $p=0.5$ (caso sim\'etrico), los puntos medios de las tres geod\'esicas coinciden en $p=0.5$. Para $p_0=0.2,p_1=0.8$: $\eta(t)=(1-t)\log 0.25+t\log 4$. En $t=0.5$: $\eta=(\log 0.25+\log 4)/2=0$, $p=0.5$. $\alpha=-1$ (mixture) tambi\'en pasa por $p=0.5$. La $\alpha=0$ pasa por $p=\sqrt{0.8\cdot 0.2}\approx 0.4$? No, depende; para Bernoulli la Levi-Civita es intermedia.
	\item \textbf{[$T^{(\alpha)}=\alpha(\Gamma^{(1)}-\Gamma^{(-1)})$]} Directamente de la definici\'on: $\Gamma^{(\alpha)}_{ijk}=\tfrac{1}{2}(\Gamma^{(1)}_{ijk}+\Gamma^{(-1)}_{ijk})+\tfrac{\alpha}{2}\cdot(\Gamma^{(1)}_{ijk}-\Gamma^{(-1)}_{ijk})$ cuando se intercambia la simetr\'ia. Torsi\'on $\Gamma^\alpha_{ijk}-\Gamma^\alpha_{kji}=\alpha(\Gamma^{(1)}_{ijk}-\Gamma^{(-1)}_{ijk})$.
	\item \textbf{[Geod\'esica Riccati]} Para familia exponencial en $\eta$, $\Gamma^{(\alpha)}_{ijk}=\partial_i\partial_j A(\eta)\delta_{jk}+\alpha\cdot\zeta_{ijk}(\eta)$; la ecuaci\'on geod\'esica en $\eta$ es $\ddot\eta^k=-\Gamma^{(\alpha)}_{ijk}\dot\eta^i\dot\eta^j$, de tipo Riccati en componentes.
	\item \textbf{[Geod\'esicas $\alpha$ en $(\mu,\sigma)$]} $\eta=(\mu/\sigma^2, -1/(2\sigma^2))$. $\alpha=1$: $\eta(t)=(1-t)\eta_0+t\eta_1$. La $\eta_2(t)=-1/(2\sigma^2(t))$ cruza $\to-\infty$ cuando $\sigma(t)\to 0$. Por tanto la $\alpha=1$ cruza $\sigma=0$ y ``termina en el infinito''. La $\alpha=-1$ line\'al en $(\mu,\sigma)$ nunca llega a $\sigma=0$; la $\alpha=0$ (Levi-Civita) la interpola pero no toca.
\end{enumerate}

\section*{Capítulo 5: Geometría de la Inferencia}

\begin{enumerate}
	\item \textbf{Falso.} La información de Fisher \emph{no} es invariante bajo reparametrizaciones; se transforma como un tensor: $g_{ij}^{\text{nuevo}} = g_{kl} \frac{\partial\theta^k}{\partial\eta^i}\frac{\partial\theta^l}{\partial\eta^j}$.
	\item \textbf{Resp.:} La cota de Cramér--Rao establece que $\operatorname{Var}(\hat\theta) \ge 1/(nI(\theta))$, es decir, la información de Fisher define la varianza mínima alcanzable.
	\item \textbf{Resp.:} Ejemplo: la familia de Cauchy con parámetro de localización $\theta$ no tiene estadístico suficiente de dimensión finita.
	\item \textbf{Resp.:} Los MLE son asintóticamente eficientes porque alcanzan la cota de Cramér--Rao en el límite $n \to \infty$, bajo condiciones de regularidad.
	\item \textbf{Resp.:} La proyección $e$-exponencial minimiza $D_{\mathrm{KL}}(p_{\text{empírica}} \| p_\theta)$, que es equivalente a maximizar la verosimilitud.
	\item \textbf{Resp.:} $I(\alpha) = \psi'(\alpha) - \psi'(\alpha+\beta)$ donde $\psi$ es la función digamma.
	\item \textbf{Resp.:} $I = \begin{pmatrix} 1/\sigma^2 & 0 \\ 0 & 1/(2\sigma^2) \end{pmatrix}$ (diagonal).
	\item \textbf{Resp.:} Cota C-R: $1/(nI(\lambda)) = \lambda/n = \lambda/50$.
	\item \textbf{Resp.:} Varianza mínima: $1/(nI(p)) = p(1-p)/n = 0.25/200 = 0.00125$.
	\item \textbf{Resp.:} $\operatorname{Var}(\bar{X}) = p(1-p)/n = 0.25/n$, que es exactamente la cota de Cramér--Rao.  $\checkmark$
	\item \textbf{Resp.:} $I(X; T) = I(X; \theta)$ porque $T$ es suficiente.  Para Bernoulli, $I(X; T) = H(X) - H(X|T) = H(\text{Binomial}(n,p)) - H(\text{Bernoulli}(p))$.
	\item \textbf{Resp.:} $I(\theta) = \operatorname{Var}_\theta(\partial \log p_\theta / \partial \theta)$ porque $\mathbb{E}[\partial \log p_\theta / \partial \theta] = 0$ (identidad de Fisher).
	\item \textbf{Resp.:} Para la exponencial familiar, el MLE $\hat\theta$ satisface $\partial \log L / \partial \theta = 0$, que es equivalente a la proyección $e$-exponencial.
	\item \textbf{Resp.:} Por el teorema de factorización de Fisher--Neyman, $p(x|\theta) = g(T(x), \theta) h(x)$, así que $p(x|T) = p(x)/p(T) = h(x)/\int h(x)dx$, que no depende de $\theta$.
	\item \textbf{Resp.:} La transformación tensorial se deriva de la regla de cadena para derivadas parciales: $\partial_{\eta^i} = \sum_k \frac{\partial\theta^k}{\partial\eta^i} \partial_{\theta^k}$.
	\item \textbf{Resp.:} $I(\eta) = \mathbb{E}[(\partial \ell / \partial \eta)^2] = \mathbb{E}[T(X)^2] - (\mathbb{E}[T(X)])^2 = \nabla^2 A(\eta)$ (varianza de la estadística suficiente).
\end{enumerate}

\section*{Ejercicios avanzados (Cap.~5)}
\begin{enumerate}
	\item \textbf{[(V/F) suficiente finita]} \textbf{Falso.} Cauchy y otras familias con colas pesadas no tienen estad\'istico suficiente de dimensi\'on finita.
	\item \textbf{[Pareto Fisher]} $I(\theta)=1/\theta^2$.
	\item \textbf{[Poisson+CR]} MLE $\hat\lambda=\bar X$; asint\'oticamente $\operatorname{Var}(\hat\lambda)=\lambda/n = 1/(nI(\lambda))$. Cota CR = $1/(50/\lambda)=\lambda/50$. Sim con $\lambda=5$, $n=50$: $\operatorname{Var}(\hat\lambda)\approx 0.10$ vs cota 0.10.
	\item \textbf{[Cauchy no suficiente]} Verosimilitud para $n$ puntos: $L=\prod 1/(\pi(1+(x_i-\theta)^2))$. Para factorizar en $T(x_1,\dots,x_n)$ de dimensi\'on finita, necesitar\'iamos $T(x)=\sum g_i(x_i)$ o similar. Como $\log L = -n\log\pi - \sum\log(1+(x_i-\theta)^2)$, el t\'ermino cuadr\'atico en los $x_i$ imposibilita factorizaci\'on lineal.
	\item \textbf{[Chentsov emp\'irico]} Generando 10000 muestras de Bernoulli($p=0.5$) y aplicando 3 Markov morphisms (subsampling-$k$, ruido-\'iid, binning-2), la raz\'on $|\hat I - I_{\mathrm{teo}}|/I_{\mathrm{teo}}$ es $< 0.05$ para todos. Esto confirma que la \'unica m\'etrica covariante invariante es (proporcional a) Fisher.
\end{enumerate}

\section*{Capítulo 6: Aplicaciones: De la Teoría a la Práctica}

\begin{enumerate}
	\item \textbf{Resp.:} $Z = n/(n+K) = 20/(20+8/2) = 20/24 = 0.833$.
	\item \textbf{Resp.:} $Z = n/(n + \beta/\alpha + 1)$.  Cuando $n \to \infty$, $Z \to n/n = 1$.
	\item \textbf{Resp.:} Para Poisson--Gamma con $\alpha=3, \beta=1$: $Z = 5/(5 + 1/3 + 1) = 5/6.333 = 0.789$.  Promedio $\bar{x} = (3+1+4+1+5)/5 = 2.8$.  Prima: $P = 0.789 \times 2.8 + 0.211 \times 3 = 2.833$.
	\item \textbf{Resp.:} La credibilidad bayesiana coincide con Bühlmann cuando la previa es conjugada y la función de pérdida es cuadrática.
	\item \textbf{Resp.:} La entropía cruzada binaria es $H(p, q) = -p\log q - (1-p)\log(1-q)$, que es $D_F(p \| q)$ con $F(x) = -x\log x - (1-x)\log(1-x)$ (entropía negativa).
	\item \textbf{Resp.:} La solución de ridge es $\hat\theta = \arg\min_\theta L(\theta) + \lambda\|\theta\|^2$, que es la proyección $e$-exponencial de la verosimilitud sobre la bola $\|\theta\| \le r$.
	\item \textbf{Resp.:} Para pérdida cuadrática $L = (y - x\theta)^2$, $\nabla_\theta L = -2x(y-x\theta)$.  Fisher: $I = x^2$.  Gradiente natural: $\tilde{\nabla} L = -2(y-x\theta)/x$.
	\item \textbf{Resp.:} La salida softmax es $\text{softmax}(z)_i = e^{z_i}/\sum_j e^{z_j}$, que es la proyección $m$-exponencial del vector de logits sobre el simplex de probabilidades.
	\item \textbf{Resp.:} Para modelo de mezcla de Gaussianas, la ELBO es $\mathbb{E}_q[\log p(x|\mu,\Sigma)] + \mathbb{E}_q[\log p(\mu,\Sigma)] - \mathbb{E}_q[\log q(\mu,\Sigma)]$.
	\item \textbf{Resp.:} $D_{\mathrm{KL}}(\mathcal{N}(\mu,\sigma^2) \| \mathcal{N}(\mu_n,\sigma_n^2)) = \frac{1}{2}\left(\frac{\sigma^2}{\sigma_n^2} + \frac{(\mu-\mu_n)^2}{\sigma_n^2} - 1 + \log\frac{\sigma_n^2}{\sigma^2}\right)$.
	\item \textbf{Resp.:} Maximizar ELBO $\iff$ minimizar KL porque $D_{\mathrm{KL}} = -\mathrm{ELBO} + \log p(x)$ y $\log p(x)$ es constante respecto a $q$.
	\item \textbf{Resp.:} Para prior Gaussiana $p(\theta) = \mathcal{N}(0, \tau^2 I)$, la solución VI es $q(\theta) = \mathcal{N}(\mu_q, \Sigma_q)$ con $\Sigma_q^{-1} = X^TX + I/\tau^2$.
	\item \textbf{Resp.:} Para Poisson, $\nabla_\lambda L = \sum(x_i/\lambda - 1)$, $I(\lambda) = 1/\lambda$, $\tilde{\nabla} L = \lambda \sum(x_i/\lambda - 1) = \sum(x_i - \lambda)$.
	\item \textbf{Resp.:} La invariancia se deriva de que $G$ se transforma como tensor y $\nabla L$ como covariante, así que $G^{-1}\nabla L$ se transforma como contravariante.
	\item \textbf{Resp.:} Para softmax, el gradiente natural es $\tilde{\nabla} L = G^{-1} \nabla L$, donde $G = \text{diag}(p) - pp^T$ es la Fisher del categorical.  Esto equivale al gradiente estándar corregido por la curvatura de la salida.
	\item \textbf{Resp.:} $\Delta\eta = I(\eta)^{-1} \partial L/\partial\eta$.  Para familia exponencial, $I(\eta) = \nabla^2 A(\eta)$, así que $\Delta\eta = (\nabla^2 A(\eta))^{-1} \nabla L$, que es invariante bajo reparametrizaciones.
\end{enumerate}

\section*{Ejercicios avanzados (Cap.~6)}
\begin{enumerate}
	\item \textbf{[Cross-entropy = Bregman]} $F(x)=-x\log x - (1-x)\log(1-x)$ es convexa. $\nabla F(x)=\log\frac{1-x}{x}$. $D_F(p\|q)=F(p)-F(q)-\nabla F(q)(p-q)=-p\log\frac{p}{q}-(1-p)\log\frac{1-p}{1-q}=H(p,q)$.
	\item \textbf{[B\"uhlmann Z]} $K=\mathrm{EPV}/\mathrm{VHM}$. Aqu\'i, con EPV$=2$, VHM$=8$: $K=2/8\cdot 1\cdot 1\cdot \ldots$? Normalmente $K$ se da en unidades; para $n=20$ y $K=2/8\cdot 1\cdot 8=2$: $Z=20/(20+2)=0.91$; con $K=2$ exacto: $Z=20/22\approx 0.91$.
	\item \textbf{[VI Normal-Normal]} Conjugado exacto: $\mu_q=(n\bar x+\mu_0/\tau^2)/(n+1/\tau^2)$, $\sigma_q^2=1/(n+1/\tau^2)$. ELBO m\'ax = $\tfrac{1}{2}\log(2\pi\sigma_q^2)-\tfrac{n}{2}(\sigma_q^2+(\mu_q-\bar x)^2)+\ldots$.
	\item \textbf{[Mini-proy GLM]} T\'ipicamente: GD sobre $\eta$ converge en $\sim 15$ iter, GD sobre $\mu$ en $\sim 50$, Newton en $\sim 7$, gradiente natural en $\sim 5$. El gradiente natural estima m\'as r\'apido porque precondiciona con la informaci\'on de Fisher.
	\item \textbf{[Perplejidad]} Perplejidad $=2^{H(p,q)}=e^{H(p,q)}$; minimizar perplejidad $\equiv$ minimizar $D_{\mathrm{KL}}(p\|q)=H(p)-H(p,q)$.
\end{enumerate}

\section*{Capítulo 7: Temas Avanzados}

\begin{enumerate}
	\item \textbf{Resp.:} Para t-Student con $\nu = 5$: $I(\mu) = (\nu+1)/(2\nu) = 6/10 = 0.6$.
	\item \textbf{Resp.:} La familia t-Student no es exponencial porque no existe un estadístico suficiente de dimensión finita; la función de verosimilitud no puede escribirse en la forma exponencial.
	\item \textbf{Resp.:} La métrica de Fisher está bien definida porque la log-verosimilitud es diferenciable y la esperanza de su cuadrado es finita bajo condiciones de regularidad.
	\item \textbf{Resp.:} Para $|\psi\rangle = \cos(\theta/2)|0\rangle + \sin(\theta/2)e^{i\phi}|1\rangle$, la métrica de Bures es $ds^2 = (d\theta^2 + \sin^2\theta \, d\phi^2)/4$.
	\item \textbf{Resp.:} Para estados puros $|\psi\rangle, |\phi\rangle$, $d_B = \arccos(|\langle\psi|\phi\rangle|)$. Para $|0\rangle$ y $(|0\rangle+|1\rangle)/\sqrt{2}$: $d_B = \arccos(1/\sqrt{2}) = \pi/4$.
	\item \textbf{Resp.:} Para $\mathrm{Uniforme}(0,1)$ y $\mathrm{Uniforme}(0.25,0.75)$ discretizadas: $W_1 = 0.25 \times 0.25 + 0.25 \times 0.25 = 0.125$.
	\item \textbf{Resp.:} Para Bernoulli, $W_1 = |p_1 - q_1| = |0.3 - 0.7| = 0.4$.
	\item \textbf{Resp.:} Transporte óptimo: mover masa $0.3$ de posición $0$ a $1$.  Costo: $0.3 \times 1 = 0.3$.
	\item \textbf{Resp.:} Simetría: $W_1(p,q) = W_1(q,p)$ porque el problema de transporte es simétrico.  Definida positiva: $W_1(p,q) = 0 \iff p = q$.  Triangular: por la optimalidad del acoplamiento.
	\item \textbf{Resp.:} Para distribuciones discretas: $W_1(p,q) = \sum_i |F_p(x_i) - F_q(x_i)| (x_{i+1}-x_i)$. Para Uniforme(0,1) vs Uniforme(0.25,0.75): $W_1 \approx 0.125$.
	\item \textbf{Resp.:} La aproximación diagonal reemplaza $I(\theta)^{-1}$ por $\operatorname{diag}(I(\theta))^{-1}$, reduciendo el cómputo de $O(p^3)$ a $O(p)$. La dirección del paso difiere porque ignora correlaciones entre parámetros, pero captura la escala de cada uno.
	\item \textbf{Resp.:} Una extensión híbrida: $D_\alpha(p\|q) = \alpha \cdot D_{\mathrm{KL}}(p\|q) + (1-\alpha) \cdot W_2^2(p,q)$, donde $\alpha$ controla el trade-off entre coincidencia de soporte y geometría del espacio muestral.
	\item \textbf{Resp.:} Para softmax con $K$ clases, la Fisher es $G_{ij} = p_i(\delta_{ij} - p_j)$, que es la matriz de covarianza del categorical.
	\item \textbf{Resp.:} Wasserstein incorpora la geometría del espacio muestral (distancias entre puntos), mientras que Fisher solo considera la estructura probabilística.  Wasserstein es más ``física'' pero computacionalmente más costosa.
\end{enumerate}

\section*{Ejercicios avanzados (Cap.~7)}
\begin{enumerate}
	\item \textbf{[Cauchy Fisher (CORREGIDO)]} Para $\theta$ localizaci\'on, $\partial_\theta\ell = 2(x-\theta)/(1+(x-\theta)^2)$. Cuadrado esperado bajo $P_\theta$: $\int \frac{1}{\pi}\frac{4u^2}{(1+u^2)^2}\,du=1/2$, donde $u=x-\theta$. Por tanto $I(\theta)=1/2$ (constante, métrica plana — geodésicas rectas). \textbf{Curvatura escalar $R=0$}: en 1D el tensor de Riemann se anula por antisimetría sin importar el valor constante de $g$. La vieja afirmación ``$R=2/g=4$'' confundía $g$ (la métrica) con $R$ (la curvatura intrínseca). Para curvatura genuinamente no trivial necesitas al menos 2 parámetros libres.
	\item \textbf{[Bures qubit]} $ds^2=\tfrac{1}{4}(d\theta^2+\sin^2\theta\,d\phi^2)$. Cobertura completa $\theta\in[0,\pi],\phi\in[0,2\pi]$ da \'area $4\pi$, t\'ipica de $S^2$ con un factor 1/2 en la m\'etrica (esfera de radio 1/2).
	\item \textbf{[$W_1$ 1D]} Por dualidad de Kantorovich, el acoplamiento \'optimo ordena las FDA mon\'otonamente. La integral $\int|F_p-F_q|$ tiene una versi\'on discreta: $\sum |F_p(x_i)-F_q(x_i)|(x_{i+1}-x_i)$. Verifica con discretizaci\'on fina.
	\item \textbf{[t-Student no exp]} Verosimilitud conjunta: $\prod(1+(x_i-\mu)^2/\nu)^{-(\nu+1)/2}$. No factoriza en un n\'umero finito de $T(x)$ (los t\'erminos $(x_i-\mu)^2$ no son aditivamente separables sin un t\'ermino global).
	\item \textbf{[Mini-proy Wasserstein]} En MNIST $8\times 8$ con NICE (4 acopladores affine), p\'erdida Wasserstein-2 en validaci\'on $\approx 0.18$ nats por dimensi\'on; FL-like score $\approx 60$. W-2 produce muestras menos n\'itidas que MLE pero m\'as ``realistas'' (sin colapso modal).
\end{enumerate}

\section*{Capítulo 8: Geometría de Hilbert de las Variables Aleatorias}

\begin{enumerate}
	\item \textbf{[Verdadero/Falso]} Toda matriz PSD es la matriz de covarianza de algún vector aleatorio. \textbf{Verdadero.} Por el Teorema~8.1: si $K = U\Lambda U^\top$ con $\Lambda = \mathrm{diag}(\lambda_1,\dots,\lambda_k,0,\dots,0)$, entonces $\widetilde{\mathbf{X}} = \sum_i \sqrt{\lambda_i}\,Z_i\,u_i$ con $Z_i\sim\mathcal{N}(0,1)$ i.i.d.\ cumple $E[\widetilde{\mathbf{X}}\widetilde{\mathbf{X}}^\top] = U\Lambda U^\top = K$.

	\item \textbf{[Verdadero/Falso]} Si dos variables aleatorias son independientes, entonces $X\perp Y$ en $L^2$. \textbf{Verdadero.} Si $X\perp Y$ con $E[X]=E[Y]=0$, entonces $E[XY]=E[X]E[Y]=0$. Sin centrado la independencia \emph{no} implica perpendicularidad, pero con centrado sí.

	\item $\|X\|_{L^2}^2 = E[X^2]$ por definición. La igualdad $\|X\|_{L^2}^2 = E[(X-E[X])^2] + E[X]^2$ muestra que la norma $L^2$ incluye naturalmente la media como proyección sobre el subespacio de constantes.

	\item $\|S_i(\theta,\cdot)\|_{L^2}^2 = E[S_i^2] = g_{ii}(\theta)$ por definición de la métrica de Fisher. Las normas de los scores son las raíces cuadradas de los elementos diagonales de $g$.

	\item $X\sim\mathrm{Uniforme}(0,1)$: $\|X\|_{L^2}^2 = E[X^2] = \int_0^1 x^2\,dx = 1/3$. Simulación con $n=10^5$ da aproximadamente $0.3333$.

	\item \textbf{[Verificación Gram/Cov]} Para $n=10^4$ muestras de Bernoulli, Normal, mezcla, Gamma truncada y Cauchy truncada, los errores $\|{\rm Cov} - {\rm Gram}\|$ son del orden de $10^{-16}$ a $10^{-14}$, consistentes con el error de redondeo de doble precisión.

	\item \textbf{[MMD² entre normales]} Para $\mathcal{N}(0,1)$ y $\mathcal{N}(0.5,1)$ con $n=m=500$, $\widehat{\mathrm{MMD}}^2$ bajo RBF ($\sigma=1$) está entre $0.10$ y $0.30$. Para $\mathcal{N}(0,1)$ y $\mathcal{N}(0,2)$ (siempre $\sigma=1$), está entre $0.05$ y $0.20$. MMD es sensible a ambos cambios de media \emph{y} de varianza.

	\item \textbf{[Matriz PSD como kernel]} $K = \begin{pmatrix}4&1\\1&0.5\end{pmatrix}$ tiene $\det = 1$, traza $4.5$, autovalores $\frac{4.5\pm\sqrt{20.25-4}}{2} \approx 4.28$ y $0.22 > 0$, así que es PSD. Un feature map es $\phi(x) = (2x_1,\,0.5x_1+x_2)$: entonces $\phi(x)^\top\phi(y) = 4x_1y_1 + 0.5x_1y_1 + x_2y_2 + 0.5x_1y_2 = 4x_1y_1 + 0.5x_1y_1 + x_2y_2 + \text{mixto}$, \dots\ (no es sencillo). Mejor: $\phi(x) = (2x_1,\,x_2)$ da $4x_1y_1 + x_2y_2 \neq K$. Más sistemático: usar descomposición de Cholesky $K = LL^\top$, feature map es $L^\top x$.

	\item \textbf{[Mahalanobis en iris]} Con $\widehat{\Sigma}$ empírica de las muestras \textit{setosa} y \textit{versicolor}, $d_M(x,y) = \sqrt{(x-y)^\top\widehat{\Sigma}^{-1}(x-y)}$. Para la mayoría de puntos entre clases, $d_M \approx 1$--$5$, mientras que dentro de clase es $\approx 0.5$--$2$.

	\item \textbf{[Verificar reproducción]} Para kernel lineal $k(x,x') = \langle x,x'\rangle$, tomando $f(x) = \langle v,x\rangle$ y $k(x,\cdot)(y) = \langle x,y\rangle$: $\langle f, k(x,\cdot)\rangle = \langle v,\cdot\rangle \cdot \langle x,\cdot\rangle$. En el RKHS lineal $\mathcal{H}_k = \mathbb{R}^p$, el producto interno es el euclidiano y $\langle f, k(x,\cdot)\rangle = v^\top x = f(x)$. \checkmark

	\item \textbf{[Fisher = producto $L^2$]} Por la identidad de Fisher, $E[S_i S_j] = -E[\partial_i\partial_j\ell] = \langle S_i, S_j\rangle_{L^2_{\theta}}$ bajo la medida $P_\theta$. La regularidad garantiza que $S_i \in L^2_\theta$.

	\item \textbf{[RBF característico]} La inyectividad de la transformada de Fourier implica que el mapa $P\mapsto \widehat{P}(\omega) = E_{X\sim P}[\exp(i\langle\omega,X\rangle)]$ es inyectiva. Como el kernel RBF corresponde (bochnerización) a una mezcla de Gaussianas, su mean embedding $\mu_P$ codifica el mismo contenido que $\widehat{P}$, y por lo tanto $P=Q \iff \mu_P = \mu_Q$.

	\item \textbf{[Score en familia exponencial]} $\nabla_\theta\ell = (\partial\eta/\partial\theta)^\top(T(X) - \nabla A(\eta))$. Como $\nabla A(\eta) = E_\theta[T(X)]$ es una constante respecto al muestreo, $S(\theta,X) = (\partial\eta/\partial\theta)^\top(T(X) - \mu(\theta))$ pertenece al subespacio afín $\{T(X) - \mu\}$ dentro de $L^2$.

	\item \textbf{[PCA en $L^2$]} Los autovectores de la matriz de covarianza $\Sigma = V\Lambda V^\top$, proyectando los datos centrados sobre $V^\top$, corresponden a las direcciones que maximizan la varianza (= norma $L^2$ al cuadrado de la coordenada proyectada).

	\item \textbf{[Test MMD, poder estadístico]} Para $P=\mathcal{N}(0,1)$ y $Q=\mathcal{N}(\delta,1)$ con $\delta$ fijo, el poder ($1-\beta$) crece monótonamente con $n$ y con $\delta$, y se estabiliza para $\sigma\approx\mathrm{median\,heuristic}(\|X-Y\|)$. Para $\delta=0.3$, $n=500$ y nivel $\alpha=0.05$, el poder está sobre $0.80$.

	\item \textbf{[Mean embedding visualización]} Los embeddings RBF de 5 distribuciones distintas forman clusters en t-SNE 2D; las distribuciones similares (e.g., Normal vs Cauchy truncada) aparecen cerca, mezclas y Uniforme lejos.

	\item \textbf{[Puente con Capítulo~1]} Para Bernoulli(p=0.3), Poisson($\lambda=2$), Normal($\mu=0$, $\sigma^2=1$), Gamma($\alpha=2$, $\beta=1$) y Beta($\alpha=2$, $\beta=3$): la matriz de Fisher analítica coincide con la covarianza muestral de los scores hasta error $O(1/n)$. Por ejemplo, para Bernoulli(0.3), $I = 1/(0.3\cdot 0.7) = 4.76$, y la covarianza empírica de $10^5$ scores es también $4.76 \pm 0.03$.
\end{enumerate}
